%%%%%%%%%%%%%%%%%%%%%%%%%%%%%%%%%%%%%%%%%%%%%%%%%%%%%%%%%%%%%%%%%%%%%%%%%
%
% Ejercicios resueltos de Análisis Matemático II.
% Doble Grado de Informática y Matemáticas.
% Universidad de Granada.
% Curso 2016/17.
% 
% 
% Colaboradores:
% Javier Sáez (@fjsaezm)
% Daniel Pozo (@danipozodg)
% Pedro Bonilla (@pedrobn23)
% Guillermo Galindo (@guillegalor)
% Antonio Coín (@antcc)
% Sofía Almeida (@SofiaAlmeida)
% Miguel Lentisco (@AlfaOmegaX)
%
% Agradecimientos:
% Andrés Herrera (@andreshp) y Mario Román (@M42) por
% las plantillas base.
%
% Sitio original:
% https://github.com/libreim/apuntesDGIIM/
%
% Licencia:
% CC BY-NC-SA 4.0 (https://creativecommons.org/licenses/by-nc-sa/4.0/)
%
%%%%%%%%%%%%%%%%%%%%%%%%%%%%%%%%%%%%%%%%%%%%%%%%%%%%%%%%%%%%%%%%%%%%%%%%%


%------------------------------------------------------------------------------
%   ACKNOWLEDGMENTS
%------------------------------------------------------------------------------

%%%%%%%%%%%%%%%%%%%%%%%%%%%%%%%%%%%%%%%%%%%%%%%%%%%%%%%%%%%%%%%%%%%%%%%%
% Plantilla básica de Latex en Español.
%
% Autor: Andrés Herrera Poyatos (https://github.com/andreshp) 
%
% Es una plantilla básica para redactar documentos. Utiliza el paquete fancyhdr
% para darle un estilo moderno pero serio.
%
% La plantilla se encuentra adaptada al español.
%
%%%%%%%%%%%%%%%%%%%%%%%%%%%%%%%%%%%%%%%%%%%%%%%%%%%%%%%%%%%%%%%%%%%%%%%%%

%%%%%%%%%%%%%%%%%%%%%%%%%%%%%%%%%%%%%%%%%%%%%%%%%%%%%%%%%%%%%%%%%%%%%%%%%
% Plantilla de Trabajo
% Modificación de una plantilla de Latex de Frits Wenneker para adaptarla 
% al castellano y a las necesidades de escribir informática y matemáticas.
%
% Editada por: Mario Román
%
% License:
% CC BY-NC-SA 3.0 (http://creativecommons.org/licenses/by-nc-sa/3.0/)
%%%%%%%%%%%%%%%%%%%%%%%%%%%%%%%%%%%%%%%%%%%%%%%%%%%%%%%%%%%%%%%%%%%%%%%%%

%%%%%%%%%%%%%%%%%%%%%%%%%%%%%%%%%%%%%%%%%%%%%%%%%%%%%%%%%%%%%%%%%%%%%%%%%
% Short Sectioned Assignment
% LaTeX Template
% Version 1.0 (5/5/12)
%
% This template has been downloaded from:
% http://www.LaTeXTemplates.com
%
% Original author:
% Frits Wenneker (http://www.howtotex.com)
%
% License:
% CC BY-NC-SA 3.0 (http://creativecommons.org/licenses/by-nc-sa/3.0/)
%
%%%%%%%%%%%%%%%%%%%%%%%%%%%%%%%%%%%%%%%%%%%%%%%%%%%%%%%%%%%%%%%%%%%%%%%%%


% Tipo de documento y opciones.
\documentclass[11pt, a4paper]{article}


%---------------------------------------------------------------------------
%   PAQUETES
%---------------------------------------------------------------------------

% Idioma y codificación para Español.
\usepackage[utf8]{inputenc}
\usepackage[spanish, es-tabla, es-lcroman, es-noquoting]{babel}
\selectlanguage{spanish} 
%\usepackage[T1]{fontenc}

% Fuente utilizada.
\usepackage{courier}    % Fuente Courier.
\usepackage{microtype}  % Mejora la letra final de cara al lector.

% Diseño de página.
\usepackage{fancyhdr}   % Utilizado para hacer títulos propios.
\usepackage{lastpage}   % Referencia a la última página.
\usepackage{extramarks} % Marcas extras. Utilizado en pie de página y cabecera.
\usepackage[parfill]{parskip}    % Crea una nueva línea entre párrafos.
\usepackage{geometry}            % Geometría de las páginas.

% Símbolos y matemáticas.
\usepackage{amssymb, amsmath, amsthm, amsfonts, amscd}
\usepackage{upgreek}

% Otros.
\usepackage{enumitem}   % Listas mejoradas.
\usepackage[hidelinks]{hyperref}
\usepackage{graphicx}   % Gráficos.


%---------------------------------------------------------------------------
%   OPCIONES PERSONALIZADAS
%---------------------------------------------------------------------------

% Redefinir letra griega épsilon.
\let\epsilon\upvarepsilon

% Formato de texto.
\linespread{1.1}            % Espaciado entre líneas.
\setlength\parindent{0pt}   % No indentar el texto por defecto.
\setlist{leftmargin=.5in}   % Indentación para las listas.

% Estilo de página.
\pagestyle{fancy}
\fancyhf{}
\geometry{left=3cm,right=3cm,top=3cm,bottom=3cm,headheight=1cm,headsep=0.5cm}   % Márgenes y cabecera.

% Ruta donde buscar gráficos
\graphicspath{{../Recursos/Plantillas/}, {Recursos/Plantillas/}}

% Redefinir entorno de demostración (reducir espacio superior)
%\makeatletter
%\renewenvironment{proof}[1][\proofname] {\vspace{-15pt}\par\pushQED{\qed}\normalfont\topsep6\p@\@plus6\p@\relax\trivlist\item[\hskip\labelsep\it#1\@addpunct{.}]\ignorespaces}{\popQED\endtrivlist\@endpefalse}
%\makeatother

% Aumentar el tamaño del interlineado
\linespread{1.3}


%---------------------------------------------------------------------------
%   COMANDOS PERSONALIZADOS
%---------------------------------------------------------------------------

% Valor absoluto: \abs{}
\providecommand{\abs}[1]{\lvert#1\rvert}    

% Fracción grande: \ddfrac{}{}
\newcommand\ddfrac[2]{\frac{\displaystyle #1}{\displaystyle #2}}

% Texto en negrita en modo matemática: \bm{}
\newcommand{\bm}[1]{\boldsymbol{#1}}

% Línea horizontal.
\newcommand{\horrule}[1]{\rule{\linewidth}{#1}}

% Letras de conjuntos
\newcommand{\R}{\mathbb{R}}
\newcommand{\N}{\mathbb{N}}

% Sucesiones
\newcommand{\xn}{\{x_n\}}
\newcommand{\fn}{\{f_n\}}


%---------------------------------------------------------------------------
%   CABECERA Y PIE DE PÁGINA
%---------------------------------------------------------------------------

% Cabecera del documento.
\renewcommand\headrule{
	\begin{minipage}{1\textwidth}
		\hrule width \hsize 
	\end{minipage}
}

% Texto de la cabecera.
\lhead{\subject}  % Izquierda.
\chead{}            % Centro.
\rhead{\docauthor}    % Derecha.

% Pie de página del documento.
\renewcommand\footrule{                                 
	\begin{minipage}{1\textwidth}
		\hrule width \hsize   
	\end{minipage}\par
}

% Texto del pie de página.
\lfoot{}                                                 % Izquierda
\cfoot{}                                                 % Centro.
\rfoot{Página\ \thepage\ de\ \protect\pageref{LastPage}} % Derecha.


%---------------------------------------------------------------------------
%   ENTORNOS PARA MATEMÁTICAS
%---------------------------------------------------------------------------

% Nuevo estilo para definiciones.
\newtheoremstyle{definition-style} % Nombre del estilo.
{10pt}               % Espacio por encima.
{10pt}               % Espacio por debajo.
{}                   % Fuente del cuerpo.
{}                   % Identación.
{\bf}                % Fuente para la cabecera.
{.}                  % Puntuación tras la cabecera.
{.5em}               % Espacio tras la cabecera.
{\thmname{#1}\thmnumber{ #2}\thmnote{ (#3)}}     % Especificación de la cabecera (actual: nombre en negrita).

% Nuevo estilo para notas.
\newtheoremstyle{remark-style} 
{10pt}                
{10pt}                
{}                   
{}                   
{\itshape}          
{.}                  
{.5em}               
{}                  

% Nuevo estilo para teoremas y proposiciones.
\newtheoremstyle{theorem-style}
{10pt}                
{10pt}                
{\itshape}           
{}                  
{\bf}             
{.}                
{.5em}               
{\thmname{#1}\thmnumber{ #2}\thmnote{ (#3)}}                   

% Nuevo estilo para ejemplos.
\newtheoremstyle{example-style}
{10pt}                
{10pt}                
{}                  
{}                   
{\scshape}              
{:}                 
{.5em}               
{}                   

% Teoremas, proposiciones y corolarios.
\theoremstyle{theorem-style}
\newtheorem{nth}{Teorema}[section]
\newtheorem*{nprop}{Proposición}
\newtheorem{ncor}{Corolario}[section]

% Definiciones.
\theoremstyle{definition-style}
\newtheorem*{ndef}{Definición}
\newtheorem{ejer}{Ejercicio}

% Notas.
\theoremstyle{remark-style}
\newtheorem*{nota}{Nota}
\newtheorem*{sol}{Solución}

% Ejemplos.
\theoremstyle{example-style}
\newtheorem{ejemplo}{Ejemplo}[section]

% Listas ordenadas con números romanos (i), (ii), etc.
\newenvironment{nlist}
{\begin{enumerate}
\renewcommand\labelenumi{(\emph{\roman{enumi})}}}
{\end{enumerate}}

% División por casos con llave a la derecha.
\newenvironment{rcases}
  {\left.\begin{aligned}}
  {\end{aligned}\right\rbrace}


%---------------------------------------------------------------------------
%   PÁGINA DE TÍTULO
%---------------------------------------------------------------------------

% Título del documento.
\newcommand{\subject}{Análisis Matemático II: Ejercicios resueltos}

% Autor del documento.
\newcommand{\docauthor}{DGIIM}

% Título
\title{
  \normalfont \normalsize 
  \textsc{Universidad de Granada} \\ [25pt]    % Texto por encima.
  \horrule{0.5pt} \\[0.4cm] % Línea horizontal fina.
  \huge Análisis Matemático II\\ \Large{Ejercicios resueltos}\\ % Título.
  \horrule{2pt} \\[0.5cm] % Línea horizontal gruesa.
}

% Autor.
\author{\Large{Doble Grado de Informática y Matemáticas}}

% Fecha.
\date{\vspace{-1.5em} \normalsize Curso 2016/17}


%---------------------------------------------------------------------------
%   COMIENZO DEL DOCUMENTO
%---------------------------------------------------------------------------

\begin{document}

\maketitle  % Título.
\vfill
\begin{center}
\includegraphics{by-nc-sa}  % Licencia.
\end{center}
\newpage


\begin{ejer}
    Probar que el espacio $\mathcal{C}_B(A,\mathbb{R}^M)$ es un espacio de Banach, esto es, un espacio normado y completo.

\begin{proof}
	Empezamos probando que $(\mathcal{C}_B(A, \mathbb R^M), \|\cdot \|_\infty)$ es un espacio normado:
	\begin{itemize}
	
\item  Positividad. En primer lugar puesto que la norma se ha definido como supremo de un conjunto de numeros positivos, tendremos que $\|f\|_\infty \geq 0$ para toda $f\in (\mathcal{C}_B(A, \mathbb R^M)$. Además, $|f\|_\infty=0 \iff sup_{x \in A}|f(x)|=0 \iff f(x)=0, \ \forall x \in A \iff f$ es la función 0.
\item Homogeneidad. Si $k \in \mathbb R \text{ entonces } \|kf\|_\infty=|k|\|f\|_\infty \iff sup_{x \in A}|kf(x)|=sup_{x \in A}|k||f(x)|=|k|sup_{x \in A}|f(x)| = |k| \| f\|_\infty$.
\item Desigualdad triangular.
$|f+g|_\infty \leq |f|_\infty + |g|_\infty\iff sup_{x \in A}|f(x)+g(x)| \leq sup_{x \in A}|f(x)| + sup_{x \in A}|g(x)|$.
para cualesquiera $f, g \in \mathcal{C}_B(A, \mathbb R^M)$.
\end{itemize}

Para demostrar que ${f_n} \to f$  c.u. en A $\iff f_n \to f\ en\ \mathcal{C}_B(A, \mathbb R^M)$, solo tenemos que observar que ${f_n} \to f$  c.u. en A  significa que: $$\forall \epsilon>0 \ \ \exists n_0\in \mathbb N \ :\ n\geq n_0 \Rightarrow |f_n(x)-f(x) |<\epsilon , \ \forall x\in A,$$
lo cual equivale a decir
$$\forall \epsilon>0 \ \ \exists n_0\in \mathbb N \ :\ n\geq n_0 \Rightarrow \|f_n-f\|_\infty = sup_{x \in A} |f_n-f| \leq \epsilon , \ \forall x\in A,$$
es decir, $f_n \to f\ en\ \mathcal{C}_B(A, \mathbb R^M)$.
 
Por último, $\mathcal{C}_B(A, \mathbb R^M)$ es de Banach si es completo, es decir si toda sucesión $\{ f_n\}$ (de funciones en $\mathcal{C}_B(A, \mathbb R^M)$) de Cauchy converge. La prueba es análoga a la que se hizo para ver que $\mathcal{C}(A,\mathbb{R}^M)$, con $A \subseteq \mathbb{R}^N$ compacto, era completo. La única diferencia será que tras probar la convergencia uniforme de $f_n$ a una función  $f$, deberemos probar que $f\in \mathcal{C}_B(A,\mathbb{R}^M)$, es decir que el límite uniforme de la sucesión $\{f_n\}$ de funciones continuas y acotadas es una función continua y acotada. Veámoslo.

Recordemos que ya sabemos por teoría que $f$ es continua. Para la acotación, tomando $\epsilon = 1$ en la definición de convergencia uniforme, obtenemos un $n_0 \in \mathbb{N}$ tal que $n \ge n_0 \Rightarrow |f_n(x) - f(x)| < 1\ \forall x \in A$. Por otro lado, como $f_{n_0}$ es acotada, existe un $M>0$ tal que $|f_{n_0}(x)| \le M\ \forall x \in A$. Entonces, se tiene que: $$|f(x)| = |f(x) - f_{n_0}(x) + f_{n_0}(x)| \le |f(x) - f_{n_0}(x)| + |f_{n_0}(x)| < 1 + M, \ \ \forall x \in A$$
Por tanto, $f$ está acotada.
\end{proof}
\end{ejer}

\begin{ejer} Probar que en el espacio $\mathcal{C}(A,\mathbb{R}^M)$, con $A \subseteq \mathbb{R}^N$ compacto, cualquier bola cerrada y centrada en el origen es homeomorfa a una bola cerrada y centrada en un punto arbitrario.

\begin{sol}
	Sea $\epsilon > 0$, y $\tilde f \in \mathcal C (A, \mathbb{R}^M)$. Consideremos la aplicación $\varphi: \overline{B}_\infty(0,\epsilon) \to \mathcal \overline{B}_\infty(\tilde f, \epsilon) $ dada por $\varphi(f) = f + \tilde f$.
	Por un lado, $\varphi$ está bien definida, pues: $$\|\tilde f - \varphi(f) \|_\infty= \| \tilde f - f - \tilde f \|_\infty = \|f\|_\infty \le \epsilon \Rightarrow \varphi(f) \in \overline{B}_\infty(\tilde f, \epsilon)$$
	
Tenemos que $\varphi$ es continua, pues la preimagen de un entorno básico (bolas abiertas) es un entorno básico. En efecto, %%% COMPLETAR

Además, la inversa de $\varphi$ existe: $\varphi^-1(g) = g - \tilde f\ \ \forall g \in \overline{B}_\infty(\tilde f, \epsilon)$. Es continua por el mismo motivo que $\varphi$.

Por tanto, $\varphi$ es un homeomorfismo.
\end{sol}	
\end{ejer}

\begin{ejer}
	
Sea $\{ f_n \}$ una sucesión de funciones reales uniformemente continuas en todo $\mathbb R$ que converge uniformemente a una función real $f$. ¿Puede concluirse que  la función $f$ es necesariamente uniformemente continua?. 

\begin{sol}
La respuesta es afirmativa. Veamos la prueba.
	
Dado $\epsilon>0$, como $f_n \rightarrow f$ converge uniformemente, $\exists k>0$
tal que $$n>k \Rightarrow |f_n(x)-f(x)| < \epsilon/3, \ \ \forall x \in A.$$
De otro lado, por ser $f_k$ uniformemente continua en $A$, 
$\exists \delta >0$ tal que $$\forall x,y \in A\ con\ |x-y| < \delta,\ se\ tiene\ |f_k(x)-f_k(y)| < \epsilon /3$$
Juntando ambas informaciones: $$|f(x)-f(y)| \leq |f(x)-f_k(x)| + |f_k(x)-f_k(y)| + |f_k(y)-f(y)| < \epsilon$$
Es decir, hemos probado que dado $\epsilon>0$, existe $\delta>0$ tal que $$\forall x,y \in A\ con\ |x-y|<\delta\ se\ verifica\ |f(x)-f(y)| < \epsilon$$ 
Por tanto, $f$ es uniformemente continua.
\end{sol}
\end{ejer}

\begin{ejer}
Estudiar la convergencia puntual y uniforme de la sucesión de funciones $f_n$ definidas en $[0,1]$ mediante $f_n(x) = x-x^n$ para todo $x\in [0,1]$.

\begin{sol}
Sabemos que para $0\le x<1$, $f_n(x)=x-x^n \to x$, y para $x=1$, tenemos que $f_n(x)=1-1^n = 0 \to 0$. Por tanto, el límite puntual es: $$f(x)= \left\{ \begin{array}{lcc} x & si & 0 \le x < 1 \\ \\ 0 & si & x=1  \end{array} \right.$$
Como cada $f_n$ es continua y $f$ no es continua, no hay convergencia uniforme. 
\end{sol}
\end{ejer}

\begin{ejer}
	Estudiad la convergencia uniforme de la sucesión de funciones $f_n$ definidas en [$0,99999]$ mediante $f_n(x) = x^n$ para todo $x\in [0,99999]$.

\begin{sol}
	En efecto, la sucesión de funciones converge uniformemente. En primer lugar, $\fn \xrightarrow {c.p} f=0$ por ser potencia de base menor que $1$. Además, por ser una función potencial, el valor máximo que toma es $0.99999^n$. Por tanto: $|x^n| \leq 0.99999^n \to 0$, luego $\fn$ converge uniformemente a $f=0$. 
\end{sol}
\end{ejer}

\begin{ejer}
	Estudiad la convergencia puntual y uniforme de la sucesión de funciones $f_n$ definidas en $[0,1]$ mediante $f_n(x) = (x-\frac{1}{n})^2$ para todo $x\in [0,1]$.
	
\begin{sol}
  Sabemos que $\{\frac{1}{n}\} \to 0$, por lo que podemos afirmar que $\{f_n(x)\} \to x^2$ puntualmente en $[0,1]$. Veamos que también hay convergencia uniforme: $$|f_n (x) - x^2| = \left|-\frac{2x}{n} + \frac{1}{n^2} \right| \leq \frac{2}{n} + \frac{1}{n^2}\to 0, \ \forall x\in [0,1].$$
\end{sol}
\end{ejer}

\begin{ejer}
	Estudiar el caracter de la siguientes series de funciones.

\item $\sum_{n\geq1} (\frac{x^n}{n!})^2 = \frac{x^{2n}}{(n!)^2}$
\item $\sum_{n\geq1} \frac{1}{x^2+n^2}$
\item $\sum_{k\geq1} \frac{x^k}{k^2}$
\item $\sum_{k\geq1} k! x^k = \sum_{k\geq1} a_k x^k$
\item $\forall x \in (-1,1) \ \ \ \ \frac{1}{(1-x)^2} = \sum_{k=1}^{\infty} kx^{k-1}$

$(i)\ \displaystyle \sum_{n\geq1} \frac{sen(nx)^2}{n^2}$

$|sen(nx)| < 1 \implies \abs{\frac{sen(nx)^2}{n^2}} \leq \frac{1}{n^2}$.

$\begin{rcases}
\abs{\frac{sen(nx)^2}{n^2}} \leq \frac{1}{n^2}\\
\sum_{n\geq1} \frac{1}{n^2}
\end{rcases}$ $\implies \sum \abs{\frac{sen^2(nx)}{n^2}} c. u \iff \sum_{n\geq1} \frac{sen^2(nx)}{n^2}$ es abs. convergente. $\implies$ \hspace*{4cm} $\sum \frac{sen^2(nx)}{n^2}$ converge uniformamente.

2.\\ 
$\begin{rcases}
\frac{x^{2n}}{(n!)^2} \leq \frac{M^{2n}}{(n!)^2} \ con \ |x| \leq M \\
\frac{M^{2n}}{(n!)^2} \ converge
\end{rcases}$
$\implies (\frac{x^{n}}{(n!)})^2$ es c.u en $[-M,M]$ 
\\\\Por que,
$lim \frac{\frac{M^{2(n+1)}}{(n+1)!)^2}}{\frac{M^{2n}}{(n!)^2}} = lim \frac{M^2}{(n+1)^2} = 0 < 1.$
\end{ejer}

\begin{ejer}
Encontrar un ejemplo de una función que cumpla:\\
$\{f_k\} \to 0$ c.u.\\
$\int_A f_k \not \to \int_A 0 = 0$\\
con $A\subset R$

Una solución es: $f_k(x) = 
\begin{cases}
1/k \ \ x < k \\
0 \ \ x \geq k

\end{cases}$
\end{ejer}
\end{document}